\section{Conclusion}
\label{sec:conclusion}

Agentic-AI safety has been treated, in the literature so far, as a property of the model. The training-time stack (RLHF, Constitutional AI), the evaluation-time stack (oversight, debate, red-teaming), and the interpretability-time stack all operate on the model as the object of analysis. Each thread has produced real progress on a real problem.

An orthogonal layer is available. Safety can additionally be a property of the substrate that admits a tool call, enforced as a typed precondition on the action's envelope rather than as a verdict on the model's output. The structural commitments of this layer are a positive TTL by construction, a typed rollback receipt slot for reversible actions, bilateral cosignature for destructive ones, and an admission predicate evaluated by a verifier external to the model. These commitments hold for mediated, correctly classified, single-envelope execution under independent cosigners, an intact registry, no bypass path, and working rollback executors.

The two layers compose. Training shapes what the model wants; admission shapes what the substrate is willing to execute. Neither layer alone is sufficient: training may fail, and admission may be too narrow if the class registry is wrong or the dispatch is bypassed. The combination is a stronger construction than either layer in isolation, and it is the construction the field would benefit from naming explicitly.

Three near-term directions follow. First, a worked benchmark of misaligned-agent behaviours against admission-layer protected and unprotected deployments, to expose the empirical gap. Second, formal-methods discipline on the class registry, so that the substrate's trust assumption is itself a checkable artifact. Third, integration with existing tool-calling protocols, so that the admission hook becomes a deployment option rather than a research prototype.

A scheming model that defeats alignment training still faces an external admission predicate before its tool call reaches an executor. In the bounded mediated path, destructive effects require destructive-class bilateral consent and reversible effects require a rollback commitment that the substrate, not the agent, checks. That is a safety property of a kind the literature has not yet had.
